\answerpaper{非数学 A 类}

\answerproblem{一、}{30}
填空题：
\begin{enumerate}[label=(\arabic*)]
  \item $\displaystyle \lim_{x\to 3}\frac{\sqrt{x^3+9}-6}{2-\sqrt{x^3-23}}=\underline{\hspace{3.5em}}$．
  \item 设 $z=f(x^2-y^2,xy)$，且 $f(u,v)$ 有连续的二阶偏导数，则 $\displaystyle \frac{\partial^2 z}{\partial x\partial y}=\underline{\hspace{3.5em}}$．
  \item 设 $f(x)=\dfrac{1}{x^2-3x+2}$，则 $f^{(n)}(0)=\underline{\hspace{3.5em}}$．
  \item 幂级数 $\displaystyle \sum_{n=1}^{\infty}\frac{(-1)^{n-1}x^{2n}}{n(2n-1)}$ 的收敛域为 $\underline{\hspace{3.5em}}$．
  \item 设曲面 $\Sigma$ 是平面 $y+z=5$ 被柱面 $x^2+y^2=25$ 所截得的部分，则
  \[
    \iint_{\Sigma}(x+y+z)\dd S=\underline{\hspace{6em}}.
  \]
\end{enumerate}
\begin{solution}
\textbf{(1)} 使用洛必达法则，得
\[
\lim_{x\to3}\frac{\sqrt{x^3+9}-6}{2-\sqrt{x^3-23}}
=\lim_{x\to3}\frac{\dfrac{3x^2}{2\sqrt{x^3+9}}}{-\dfrac{3x^2}{2\sqrt{x^3-23}}}
=-\lim_{x\to3}\frac{\sqrt{x^3-23}}{\sqrt{x^3+9}}
=-\frac{1}{3}.
\]

\textbf{(2)}
\begin{align*}
z_x&=2xf_1+yf_2,\\
z_{xy}&=2x\bigl(f_{11}(-2y)+f_{12}x\bigr)+f_2+y\bigl(f_{21}(-2y)+f_{22}x\bigr)\\
&=f_2-4xyf_{11}+2(x^2-y^2)f_{12}+xyf_{22}.
\end{align*}

\textbf{(3)}
\begin{align*}
f(x)&=-\frac{1}{x-1}+\frac{1}{x-2},\\
f^{(n)}(x)&=(-1)^{n+1}n!\left(\frac{1}{(x-1)^{n+1}}-\frac{1}{(x-2)^{n+1}}\right),\\
f^{(n)}(0)&=n!\left(1-\frac{1}{2^{n+1}}\right).
\end{align*}

\textbf{(4)} 因为 $\displaystyle \lim_{n\to\infty}\sqrt[n]{\frac{1}{n(2n-1)}}=1$，所以收敛半径为 $1$．

当 $x=\pm1$ 时，$\displaystyle \sum_{n=1}^{\infty}\frac{(-1)^{n-1}}{n(2n-1)}$ 绝对收敛，故收敛域为 $[-1,1]$．

\textbf{(5)} $\Sigma$ 的方程为 $z=5-y$，故
\[
\dd S=\sqrt{2}\dd x\dd y.
\]
$\Sigma$ 在 $xOy$ 平面的投影 $D_{xy}:x^2+y^2\le25$，故
\[
I=\sqrt{2}\iint_{D_{xy}}(x+5)\dd x\dd y
=5\sqrt{2}\iint_{D_{xy}}\dd x\dd y
=125\sqrt{2}\pi.
\]
\end{solution}

\answerproblem{二、}{14}
解方程
\[
(x^2+y^2+3)\frac{\dd y}{\dd x}=2x\left(2y-\frac{x^2}{y}\right).
\]
\begin{solution}
原方程变形为
\[
\frac{y\dd y}{x\dd x}=\frac{2(2y^2-x^2)}{x^2+y^2+3}.
\]
令 $u=x^2$，$v=y^2$，则原方程化为
\[
\frac{\dd v}{\dd u}=\frac{2(2v-u)}{u+v+3}.
\]
解方程 $2v-u=0$，$u+v+3=0$，得到 $u=-2$，$v=-1$，再令 $U=u+2$，$V=v+1$，上述方程化为
\[
\frac{\dd V}{\dd U}=\frac{2(2V-U)}{U+V}.
\]
作变量替换 $W=\frac{V}{U}$ 得到
\[
U\frac{\dd W}{\dd U}=-\frac{W^2-3W+2}{W+1}.
\]
这是分离变量方程，解之得
\[
U(W-2)^3=C(W-1)^2,
\]
回代得
\[
(y^2-2x^2-3)^3=C(y^2-x^2-1)^2.
\]
\end{solution}

\answerproblem{三、}{14}
设 $\Sigma_1$ 是以 $(0,4,0)$ 为顶点且与曲面
\[
\Sigma_2:\ \frac{x^2}{3}+\frac{y^2}{4}+\frac{z^2}{3}=1\quad(y>0)
\]
相切的圆锥面，求曲面 $\Sigma_1$ 与 $\Sigma_2$ 所围成的空间区域的体积．
\begin{solution}
设 $L$ 是 $xOy$ 平面上过点 $(0,4)$ 且与 $\frac{x^2}{3}+\frac{y^2}{4}=1$ 相切于点 $(x_0,y_0)$ 的直线，则
\[
\frac{x_0^2}{3}+\frac{y_0^2}{4}=1
\quad\text{且切线斜率}\quad
\frac{y_0-4}{x_0}=-\frac{4x_0}{3y_0},
\]
解得 $x_0=\pm\frac{3}{2}$，$y_0=1$．

显然，$\Sigma_1$ 与 $\Sigma_2$ 分别是切线 $L$ 和曲线 $\frac{x^2}{3}+\frac{y^2}{4}=1$ 绕 $y$ 轴旋转而成的曲面，它们的交线位于平面 $y=1$ 上．

记该平面与 $\Sigma_1$、$\Sigma_2$ 围成的空间区域分别记为 $\Omega_1$ 和 $\Omega_2$．

由于 $\Sigma_1$ 是底面圆的半径为 $\frac{3}{2}$，高为 $3$ 的圆锥体，所以其体积
\[
V_1=\frac{1}{3}\pi\left(\frac{3}{2}\right)^2\cdot3=\frac{9\pi}{4}.
\]
又 $\Omega_2$ 的体积为
\[
V_2=\iiint_{\Omega_2}\dd V
=\int_1^2\dd y\iint_{x^2+z^2\le3\left(1-\frac{y^2}{4}\right)}\dd x\dd z
=\pi\int_1^2 3\left(1-\frac{y^2}{4}\right)\dd y
=\frac{5\pi}{4}.
\]
因此，曲面 $\Sigma_1$ 与 $\Sigma_2$ 所围成的空间区域的体积为
\[
\frac{9\pi}{4}-\frac{5\pi}{4}=\pi.
\]
\end{solution}

\answerproblem{四、}{14}
设
\[
I_n=n\int_1^a\frac{\dd x}{1+x^n},
\]
其中 $a>1$．求极限 $\displaystyle \lim_{n\to\infty}I_n$．
\begin{solution}
记 $b=\frac{1}{a}$，则 $0<b<1$．作变量替换 $x=\frac{1}{t}$，得到
\[
I_n=\int_b^1\frac{nt^{n-1}}{t(1+t^n)}\dd t
=\int_b^1\frac{\dd\left(\ln(1+t^n)\right)}{t}.
\]
分部积分得
\[
I_n=\ln2-\frac{\ln(1+b^n)}{b}+\int_b^1\frac{\ln(1+t^n)}{t^2}\dd t.
\]
当 $t\in[b,1]$ 时，$\displaystyle \frac{\ln(1+t^n)}{t^2}\le t^{n-2}$，
\[
0\le\int_b^1\frac{\ln(1+t^n)}{t^2}\dd t
\le\int_b^1 t^{n-2}\dd t
=\frac{1-b^{n-1}}{n-1}.
\]
显然，$\displaystyle \lim_{n\to\infty}\frac{1-b^{n-1}}{n-1}=0$．由夹逼准则，$\displaystyle \lim_{n\to\infty}\int_b^1\frac{\ln(1+t^n)}{t^2}\dd t=0$．

又 $\displaystyle \lim_{n\to\infty}\frac{\ln(1+b^n)}{b}=0$，$\displaystyle \lim_{n\to\infty}I_n=\ln2$．
\end{solution}

\answerproblem{五、}{14}
设 $f(x)$ 在 $[0,1]$ 上有连续的导数且 $f(0)=0$．求证：
\[
\int_0^1 f^2(x)\dd x\le4\int_0^1(1-x)^2|f'(x)|^2\dd x,
\]
并求使上式成为等式的 $f(x)$．
\begin{solution}
由分部积分法
\begin{align*}
\int_0^1 f^2(x)\dd x
&=\left.(x-1)f^2(x)\right|_0^1-\int_0^1(x-1)2f(x)f'(x)\dd x\\
&=2\int_0^1(1-x)f'(x)\cdot f(x)\dd x.
\end{align*}
由 Cauchy 积分不等式，有
\[
\int_0^1(1-x)f'(x)\cdot f(x)\dd x
\le\left(\int_0^1(1-x)^2(f'(x))^2\dd x\right)^{\frac{1}{2}}
\left(\int_0^1 f^2(x)\dd x\right)^{\frac{1}{2}}.
\]
于是
\[
\int_0^1 f^2(x)\dd x\le4\int_0^1(1-x)^2|f'(x)|^2\dd x.
\]
等式成立时应有常数 $c$ 使得 $(1-x)f'(x)=cf(x)$．因此当 $x\in(0,1)$ 时，有
\[
\left((1-x)^cf(x)\right)'
=(1-x)^{c-1}\left((1-x)f'(x)-cf(x)\right)=0.
\]
因而存在常数 $d$ 使得 $f(x)=d(1-x)^{-c}$（$0<x<1$）．
当 $x\to0$ 时，$f(x)\to0$，故 $d=0$．于是 $f\equiv0$．所以使得题中不等式成为等式的函数是 $f(x)=0$．
\end{solution}

\answerproblem{六、}{14}
设数列 $\{x_n\}$ 满足
\[
x_0=\frac{1}{3},\qquad x_{n+1}=\frac{x_n^2}{1-x_n+x_n^2},\quad n\ge0.
\]
证明：无穷级数 $\displaystyle \sum_{n=0}^{\infty}x_n$ 收敛并求其和．
\begin{solution}
\textbf{方法 1.} 根据数学归纳法可知 $x_n>0$．

此外，
\[
x_{n+1}-x_n=-\frac{x_n(1-x_n)^2}{1-x_n+x_n^2}<0.
\]
故 $x_n$ 单调递减，$x_n\le\frac{1}{3}$．

于是，
\[
x_{n+1}=x_n\cdot\frac{x_n}{1-x_n+x_n^2}\le\frac{4}{9}x_n,
\]
$x_n$ 收敛于 $0$．

令 $f(x)=\frac{x}{1+x}$，$x>0$，不难验证 $f(x)$ 严格单调递增且其反函数为 $f^{-1}(x)=\frac{x}{1-x}$．
注意到 $x_{n+1}=f\left(f^{-1}(x_n)-x_n\right)$，故
\[
f^{-1}(x_{n+1})=f^{-1}(x_n)-x_n,\qquad
x_n=f^{-1}(x_n)-f^{-1}(x_{n+1}).
\]
\[
\sum_{i=0}^n x_i=f^{-1}(x_0)-f^{-1}(x_{n+1}).
\]
\[
\sum_{i=0}^{\infty}x_i=f^{-1}(x_0)-f^{-1}(0)=\frac{1}{2}.
\]

\textbf{方法 2.} 证明 $x_n$ 收敛于 $0$ 同方法 1．

注意到
\[
x_n=\frac{x_n}{1-x_n}-\frac{x_{n+1}}{1-x_{n+1}}.
\]
\[
\sum_{i=0}^n x_i=\frac{x_0}{1-x_0}-\frac{x_{n+1}}{1-x_{n+1}}.
\]
\[
\sum_{i=0}^{\infty}x_i=\frac{x_0}{1-x_0}=\frac{1}{2}.
\]
\end{solution}
